% ============================================================
% Préambule — PT Foundations (article FR, 40 p)
% Adapté de PT_MONOGRAPHY/preamble.tex
% ============================================================

\documentclass[11pt,a4paper]{article}

% --- Encodage et langue ---
\usepackage[utf8]{inputenc}
\usepackage[T1]{fontenc}
\usepackage[french]{babel}
\usepackage{csquotes}

% --- Mathématiques ---
\usepackage{amsmath,amssymb,amsthm,mathtools}
\usepackage{bm}

% --- Mise en page ---
\usepackage[margin=2.5cm]{geometry}
\usepackage{microtype}
\usepackage{enumitem}

% --- Tableaux et figures ---
\usepackage{booktabs}
\usepackage{graphicx}
\usepackage{tikz}
\usetikzlibrary{arrows.meta,positioning,fit,calc}

% --- Bibliographie ---
\usepackage[backend=biber,style=numeric,sorting=none]{biblatex}
\addbibresource{references.bib}

% --- Hyperliens ---
\usepackage[colorlinks=true,linkcolor=blue!60!black,citecolor=green!50!black,urlcolor=blue!50!black]{hyperref}
\usepackage{cleveref}

% --- Environnements de théorème ---
\theoremstyle{plain}
\newtheorem{theoreme}{Théorème}[section]
\newtheorem{proposition}[theoreme]{Proposition}
\newtheorem{lemme}[theoreme]{Lemme}
\newtheorem{corollaire}[theoreme]{Corollaire}

\theoremstyle{definition}
\newtheorem{definition}[theoreme]{Définition}
\newtheorem{convention}[theoreme]{Convention}

\theoremstyle{remark}
\newtheorem{remarque}[theoreme]{Remarque}
\newtheorem{exemple}[theoreme]{Exemple}
\newtheorem*{conjecture*}{Conjecture}
\newtheorem{conjecture}[theoreme]{Conjecture}

% --- Macros PT canoniques ---
% Secteurs (convention 2026 : q_+ / q_-, jamais q_stat/q_therm)
\newcommand{\qplus}{q_{+}}
\newcommand{\qminus}{q_{-}}

% Constantes structurelles
\newcommand{\mustar}{\mu^{*}}
\newcommand{\sper}{s}

% Dimensions anomales
\newcommand{\gammap}[1]{\gamma_{#1}}
\newcommand{\gammathree}{\gamma_{3}}
\newcommand{\gammafive}{\gamma_{5}}
\newcommand{\gammaseven}{\gamma_{7}}

% Phases cycliques
\newcommand{\thetap}[1]{\theta_{#1}}
\newcommand{\sinp}[1]{\sin^{2}\theta_{#1}}

% Status tags (10/10 [THM])
\newcommand{\THM}{\textsc{[Thm]}}

% --- Théorèmes T-series (nommés selon NOMENCLATURE_MAP.md) ---
\newcommand{\Tzero}{\textbf{T0}}
\newcommand{\Tone}{\textbf{T1}}
\newcommand{\Ttwo}{\textbf{T2}}
\newcommand{\Tthree}{\textbf{T3}}
\newcommand{\Tfour}{\textbf{T4}}
\newcommand{\Tfive}{\textbf{T5}}
\newcommand{\Tsix}{\textbf{T6}}

% --- Axiomes de pont BA-series (théorèmes prouvés) ---
\newcommand{\BAzero}{\textbf{BA0}}
\newcommand{\BAone}{\textbf{BA1}}
\newcommand{\BAtwo}{\textbf{BA2}}
\newcommand{\BAthree}{\textbf{BA3}}
\newcommand{\BAfour}{\textbf{BA4}}
\newcommand{\BAfive}{\textbf{BA5}}

% --- Lemme L0 (entropie maximale) ---
\newcommand{\Lzero}{\textbf{L0}}

% --- Macros sp\'ecifiques Article 2 ---
\newcommand{\Cper}{\mathcal{C}_{\rm per}}
\newcommand{\genus}{g}
\newcommand{\gPT}{g_{\rm PT}}
\newcommand{\Ric}{\mathrm{Ric}}
\newcommand{\Hess}{\mathrm{Hess}}
\newcommand{\Rres}{R}
\newcommand{\HBK}{H_{\rm BK}}
\newcommand{\lambdaH}{\lambda_H}
\newcommand{\DeltaMaslov}{\Delta_N^{\rm Maslov}}

% --- M\'etadonn\'ees article ---
\title{La Th\'eorie de la Persistance : lectures g\'eom\'etrique et spectrale}
\author{Yan Senez}
\date{\today}
